% META: {"id":"1SPE-SUITES-CO-042","chapitre":"1SPE-SUITES","type_objet":"corrige","exercice_ref":"1SPE-SUITES-EX-042","capacites":["1SPE-SUITES-C4","1SPE-SUITES-C5","1SPE-SUITES-C8"],"capacites_codes":["C4","C5","C8"],"sources_inspiration":[],"status":"approved","origin":"generated","status_history":["generated","audited","accepted","approved"],"release_acceptance":"RELEASE_OWNER_BATCH_ACCEPTANCE","acceptance_closure_digest":"sha256:9b3ccf9a81c5520fbb7e03b7b2d3e7bf2057a02834b6d908bf3be5f49f3b553f","acceptance_receipt_digest":"sha256:4fad86f0282e0fa4e0419cca1ad6379ae7af5a280c04a4a90880f90a1c044242"}
% BEGIN-VERIFY
% from fractions import Fraction
% u = lambda k: 2 * Fraction(3, 4)**k
% S = lambda m: sum((u(k) for k in range(m + 1)), Fraction(0))
% formula = lambda m: 8 * (1 - Fraction(3, 4)**(m + 1))
% for m in (0, 1, 2, 10, 30, 100):
%     assert S(m) == formula(m)
%     assert S(m) < 8
% v = lambda k: 3 * Fraction(1, 2)**k
% T = lambda m: sum((v(k) for k in range(m + 1)), Fraction(0))
% formula_T = lambda m: 6 * (1 - Fraction(1, 2)**(m + 1))
% for m in (0, 1, 2, 10):
%     assert T(m) == formula_T(m)
% END-VERIFY
\begin{corrige}{1SPE-SUITES-EX-042}

\textbf{1. Premiers termes et premières sommes.}

\[
u_0=2,\qquad u_1=\frac32,\qquad u_2=\frac98.
\]
\[
S_0=2,\qquad S_1=2+\frac32=\frac72,\qquad
S_2=\frac72+\frac98=\frac{37}{8}.
\]
Ces valeurs augmentent tout en semblant se rapprocher d'une valeur finie.

\textbf{2. Expression de $S_n$.}

Les termes additionnés sont ceux d'une suite géométrique de premier terme $2$
et de raison $\frac34$. Comme $\frac34\neq1$, la formule de somme donne
\[
S_n=2\,\frac{1-\left(\frac34\right)^{n+1}}{1-\frac34}
=8\left(1-\left(\frac34\right)^{n+1}\right).
\]

\textbf{3. Variation et majoration.}

Par définition des sommes,
\[
S_{n+1}-S_n=u_{n+1}>0.
\]
La suite $(S_n)$ est donc strictement croissante. De plus,
$\left(\frac34\right)^{n+1}>0$, donc
\[
S_n=8\left(1-\left(\frac34\right)^{n+1}\right)<8.
\]
Les sommes augmentent à chaque nouveau terme, tout en restant inférieures à
$8$.

\textbf{4. Observation numérique.}

\[
S_{10}\approx7{,}662,\qquad
S_{30}\approx7{,}999,\qquad
S_{100}\approx8{,}000.
\]
Les calculs suggèrent que les valeurs se rapprochent de $8$. On conjecture ce
comportement à partir des observations numériques, \textbf{sans preuve formelle
de limite}.

\textbf{5. Transfert à une autre somme géométrique.}

Les termes $(v_n)$ forment une suite géométrique de premier terme $3$ et de
raison $\frac12$. Ainsi
\[
T_n=3\,\frac{1-\left(\frac12\right)^{n+1}}{1-\frac12}
=6\left(1-\left(\frac12\right)^{n+1}\right).
\]
Comme $\left(\frac12\right)^{n+1}>0$, on a $T_n<6$. Au rang $10$,
\[
T_{10}=6\left(1-\left(\frac12\right)^{11}\right)
\approx5{,}997.
\]
On compare les écarts aux valeurs observées :
\[
8-S_{10}\approx0{,}338
\qquad\text{et}\qquad
6-T_{10}\approx0{,}003.
\]
Au rang $10$, $T_{10}$ est donc plus proche de $6$ que $S_{10}$ ne l'est de
$8$.

\end{corrige}
